% Plan 2 WIP bridge block.
% Purpose: replacement import for the Fusco--Julin centre and radial
% trace portion of WIP_WeightedMetricTrace.tex.
%
% This file is intentionally a TeX fragment, not a standalone document:
% it assumes the importing file has loaded amsmath/amsthm/amssymb and has
% theorem, lemma, proposition, and remark environments.

\providecommand{\R}{\mathbb{R}}
\providecommand{\rstar}{\rho_*}
\providecommand{\Hn}{\mathcal H^{n-1}}

\section{Same-centre Fusco--Julin bridge}
\label{sec:fjbridge}

This block isolates the exact centre package needed by the weighted
metric trace argument.  It deliberately does not replace a
Fusco--Julin centre by the bulk centroid.  All estimates below use the
same centre \(z_\rho\) for the centred symmetric difference, the normal
oscillation, and the oriented radial trace.

\subsection{Normalised notation}
\label{ssec:fjbridge-notation}

Let \(E\subset\R^n\) be a set of finite perimeter with
\(0<|E|<\infty\), and write
\[
   R_E:=\left(\frac{|E|}{\omega_n}\right)^{1/n},
   \qquad B_E(z):=B_{R_E}(z),
   \qquad P_E^*:=P(B_{R_E})=n\omega_n R_E^{n-1}.
\]
The isoperimetric deficit is
\[
   \mathcal D(E):=\frac{P(E)-P_E^*}{P_E^*}.
\]
For \(z\in\R^n\) define the centred asymmetry and the pointed
Fusco--Julin oscillation by
\[
   \alpha_{\rm FJ}(E,z)
   :=\frac{|E\Delta B_E(z)|}{|B_{R_E}|},
\]
\[
   \beta_{\rm FJ}(E,z)^2
   :=\frac12\int_{\partial^*E}
      \left|\nu_E(x)-\frac{x-z}{|x-z|}\right|^2\,d\Hn(x),
\]
with any value assigned to the radial vector at \(x=z\).  Since
\(\{z\}\) is \(\Hn\)-null for \(n\ge2\), this convention is irrelevant.
The scale-invariant version is
\[
   b_{\rm FJ}(E,z)^2:=\frac{\beta_{\rm FJ}(E,z)^2}{P_E^*}
   =\frac{1}{2P_E^*}\int_{\partial^*E}
      \left|\nu_E(x)-\frac{x-z}{|x-z|}\right|^2\,d\Hn(x).
\]

\begin{remark}[No collision with the coarea symbol]
The symbol \(\beta_{\rm FJ}\) in this block is only the
Fusco--Julin normal oscillation.  It is not the level-set flux
\(\beta(t)=\int_{\partial^*E_t}|\nabla u|\,d\Hn\) used in the
level-set identity.
\end{remark}

\begin{lemma}[Scaling of the pointed quantities]
\label{lem:fjbridge-scaling}
Let \(\lambda>0\), \(E_\lambda:=\lambda E\), and
\(z_\lambda:=\lambda z\).  Then
\[
   \alpha_{\rm FJ}(E_\lambda,z_\lambda)
      =\alpha_{\rm FJ}(E,z),
   \qquad
   b_{\rm FJ}(E_\lambda,z_\lambda)
      =b_{\rm FJ}(E,z),
\]
and
\[
   \beta_{\rm FJ}(E_\lambda,z_\lambda)^2
      =\lambda^{n-1}\beta_{\rm FJ}(E,z)^2,
   \qquad
   \mathcal D(E_\lambda)=\mathcal D(E).
\]
\end{lemma}

\begin{proof}
Volumes scale by \(\lambda^n\), perimeters and reduced-boundary
measures by \(\lambda^{n-1}\), and the measure-theoretic normal is
invariant under dilation.  The radial direction is also invariant:
\[
   \frac{\lambda x-\lambda z}{|\lambda x-\lambda z|}
   =\frac{x-z}{|x-z|}.
\]
The displayed identities follow directly.
\end{proof}

\subsection{The exact same-centre input}
\label{ssec:fjbridge-same-centre-input}

\begin{theorem}[Same-centre Fusco--Julin input]
\label{thm:fjbridge-same-centre-input}
For every \(n\ge2\) there is \(C_{\rm scFJ}(n)\) such that every
set \(E\subset\R^n\) of finite perimeter with \(0<|E|<\infty\)
satisfies
\[
   \inf_{z\in\R^n}
   \left[
      \alpha_{\rm FJ}(E,z)^2+b_{\rm FJ}(E,z)^2
   \right]
   \le C_{\rm scFJ}(n)\,\mathcal D(E).
\]
Equivalently, there is one centre \(z_E\) for which both the centred
symmetric-difference term and the pointed normal-oscillation term are
controlled by the same deficit.
\end{theorem}

\begin{remark}[Status of the input]
\label{rem:fjbridge-input-status}
The weighted metric trace requires Theorem
\ref{thm:fjbridge-same-centre-input} in this same-centre form.  A
statement controlling
\[
   \inf_z \alpha_{\rm FJ}(E,z)^2+\inf_z b_{\rm FJ}(E,z)^2
\]
is not sufficient, because the radial trace uses one fixed centre in
both terms.  Thus, if the available citation to Fusco--Julin
\cite[Thm.~1.1]{FJ2014} is only an unpointed or separately-minimised
form, the missing same-centre formulation remains an external
dependency.  The rest of this block is conditional on Theorem
\ref{thm:fjbridge-same-centre-input}.
\end{remark}

\begin{corollary}[Scaled same-centre package]
\label{cor:fjbridge-scaled}
Fix \(\rstar\in(0,1)\).  Let \(\rho\in[\rstar,1]\), and let
\(E\subset\R^n\) be a set of finite perimeter with
\(|E|=\omega_n\rho^n\).  Then there exists \(z\in\R^n\) such that
\[
   \left(\frac{|E\Delta B_\rho(z)|}{\omega_n\rho^n}\right)^2
   +\frac{1}{2P(B_\rho)}
      \int_{\partial^*E}
      \left|\nu_E(x)-\frac{x-z}{|x-z|}\right|^2\,d\Hn(x)
   \le C_{\rm scFJ}(n)\,\mathcal D(E).
\]
Consequently,
\[
   |E\Delta B_\rho(z)|^2
      \le C(n,\rstar)\,
         \bigl(P(E)-P(B_\rho)\bigr),
\]
and
\[
   \int_{\partial^*E}
      \left|\nu_E(x)-\frac{x-z}{|x-z|}\right|^2\,d\Hn(x)
      \le C(n)\,\bigl(P(E)-P(B_\rho)\bigr).
\]
\end{corollary}

\begin{proof}
Apply Theorem~\ref{thm:fjbridge-same-centre-input} to
\(\rho^{-1}E\) and scale back with Lemma~\ref{lem:fjbridge-scaling}.
The last two estimates use
\(\mathcal D(E)=(P(E)-P(B_\rho))/P(B_\rho)\),
\(P(B_\rho)=n\omega_n\rho^{n-1}\), and \(\rho\ge\rstar\).
\end{proof}

\subsection{Selection and localisation on good levels}
\label{ssec:fjbridge-selection}

In the torsion application let \(E_\rho=\{u>t(\rho)\}\subset B_R\),
\(|E_\rho|=\omega_n\rho^n\), and
\[
   \mathcal I(\rho):=D_I(t(\rho))
   =(P(E_\rho)-P(B_\rho))(P(E_\rho)+P(B_\rho)).
\]
For a.e. finite-perimeter level,
\[
   P(E_\rho)-P(B_\rho)
   \le \frac{\mathcal I(\rho)}{P(B_\rho)}
   \le C(n,\rstar)\mathcal I(\rho).
\]

\begin{proposition}[Good-level same-centre estimates]
\label{prop:fjbridge-good-level}
Assume Theorem~\ref{thm:fjbridge-same-centre-input}.  There are
\(\eta_I=\eta_I(n,\rstar)>0\) and \(C=C(n,R,\rstar)\) such that for
every a.e. finite-perimeter level \(\rho\in[\rstar,1]\) with
\(\mathcal I(\rho)\le\eta_I\), one can choose a centre \(z_\rho\)
with
\[
   |z_\rho|\le R+1,
\]
\[
   |E_\rho\Delta B_\rho(z_\rho)|^2
      \le C\,\mathcal I(\rho),
\]
and
\[
   \int_{\partial^*E_\rho}
      \left|\nu_\rho(x)-\frac{x-z_\rho}{|x-z_\rho|}\right|^2
      d\Hn(x)
      \le C\,\mathcal I(\rho).
\]
\end{proposition}

\begin{proof}
Corollary~\ref{cor:fjbridge-scaled} gives a same-centre \(z_\rho\).
Combining its two unnormalised estimates with the preceding bound on
\(P(E_\rho)-P(B_\rho)\) gives the two displayed quantitative bounds.

It remains to localise the centre.  Let
\[
   A_\rho:=
   \frac{|E_\rho\Delta B_\rho(z_\rho)|}{|B_\rho|}.
\]
The first quantitative bound gives \(A_\rho\le C(n,\rstar)
\mathcal I(\rho)^{1/2}\).  Choose \(\eta_I\) so small that
\(A_\rho<2\).  Then \(E_\rho\cap B_\rho(z_\rho)\) has positive measure.
Since \(E_\rho\subset B_R\), the two balls \(B_R\) and
\(B_\rho(z_\rho)\) intersect, and therefore
\[
   |z_\rho|\le R+\rho\le R+1.
\]
\end{proof}

\begin{proposition}[Measurable same-centre selection]
\label{prop:fjbridge-measurable-selection}
Let
\[
   G_I:=\{\rho\in[\rstar,1]:\mathcal I(\rho)\le\eta_I\}
\]
and \(K:=\overline{B_{R+1}}\).  Suppose the level family has the
standard Borel trace parametrisation: the map
\[
   (\rho,z)\mapsto
   \Phi(\rho,z):=
   \alpha_{\rm FJ}(E_\rho,z)^2+b_{\rm FJ}(E_\rho,z)^2
\]
is Borel on the a.e. finite-perimeter levels and continuous in \(z\)
for each such \(\rho\).  Then there is a Borel map
\(\rho\mapsto z_\rho\), with \(z_\rho\in K\) on \(G_I\), such that the
three conclusions of Proposition~\ref{prop:fjbridge-good-level} hold
for a.e. \(\rho\in G_I\).  Outside \(G_I\) one may set \(z_\rho=0\).
\end{proposition}

\begin{proof}
For \(\rho\in G_I\), Proposition~\ref{prop:fjbridge-good-level}
ensures that the same-centre admissible set inside \(K\) is non-empty.
Fix a constant \(C_0>C_{\rm scFJ}(n)\) large enough to absorb the
normalisations in Corollary~\ref{cor:fjbridge-scaled}, and define
\[
   \Gamma(\rho):=
   \{z\in K:\Phi(\rho,z)\le C_0\mathcal D(E_\rho)\}
   \qquad(\rho\in G_I).
\]
The sets \(\Gamma(\rho)\) are non-empty compact subsets of \(K\).  Their
graph is Borel because \(\Phi\) and \(\rho\mapsto\mathcal D(E_\rho)\)
are Borel.  The Kuratowski--Ryll-Nardzewski measurable selection
theorem, in the Castaing--Valadier form \cite[Thm.~III.6]{CastaingValadier1977},
therefore gives a Borel selector.  The estimates are exactly those of
Proposition~\ref{prop:fjbridge-good-level}.
\end{proof}

\begin{remark}[Trace-parametrisation dependency]
\label{rem:fjbridge-trace-param-dependency}
The only measurability input not proved in this block is the Borel
trace parametrisation of \(\Phi\).  For torsion superlevel sets this
should be supplied from the coarea disintegration of
\(\Hn\lfloor\partial^*E_\rho\) and the Borel representative of the
measure-theoretic normal.  If that parametrisation is not already in
the preliminaries, it is a separate GMT dependency; it is independent
of the same-centre issue.
\end{remark}

\subsection{Oriented radial trace}
\label{ssec:fjbridge-radial-trace}

\begin{theorem}[Finite-perimeter oriented radial trace]
\label{thm:fjbridge-radial-trace}
Let \(n\ge2\), \(R\ge1\), \(\rstar\in(0,1)\), and
\(\rho\in[\rstar,1]\).  Let \(E\subset B_R\) be a set of finite
perimeter with \(|E|=|B_\rho|\), and let \(z\in\overline{B_{R+1}}\).
Set
\[
   e_z(x):=\frac{x-z}{|x-z|},
   \qquad
   w_z(x):=\left|\frac{|x-z|}{\rho}-1\right|,
   \qquad
   X_z(x):=w_z(x)e_z(x),
\]
away from \(x=z\), with arbitrary values at \(x=z\).  Then
\[
   \int_{\partial^*E}w_z\,d\Hn
   \le
   \frac{n}{\rstar}|E\Delta B_\rho(z)|
   +\frac{W(R,\rstar)}{2}
      \int_{\partial^*E}|\nu_E-e_z|^2\,d\Hn,
\]
where one may take
\[
   W(R,\rstar):=1+\frac{2R+1}{\rstar}.
\]
No centroid containment, ray slicing, or smoothness of
\(\partial E\) is used.
\end{theorem}

\begin{proof}
By translating, take \(z=0\), write \(B=B_\rho\),
\(e=x/|x|\), \(w(r)=|r/\rho-1|\), and \(X=w(|x|)e\).  Since
\(E\subset B_R\) and \(|z|\le R+1\) before the translation, on
\(\partial^*E\) one has \(w\le W(R,\rstar)\).

On \(\partial^*E\), outside the \(\Hn\)-null point \(\{0\}\),
\[
   w=w\,e\cdot\nu_E+w(1-e\cdot\nu_E).
\]
Thus
\[
   \int_{\partial^*E}w\,d\Hn
   =
   \int_{\partial^*E}X\cdot\nu_E\,d\Hn
   +\frac12\int_{\partial^*E}w|\nu_E-e|^2\,d\Hn.
\]
The second term is at most
\[
   \frac{W(R,\rstar)}{2}
   \int_{\partial^*E}|\nu_E-e|^2\,d\Hn.
\]

It remains to estimate the oriented term.  For
\(\varepsilon>0\) let \(\eta_\varepsilon\) be a smooth radial cut-off
which is \(0\) on \(B_{\varepsilon/2}\), \(1\) outside \(B_\varepsilon\),
and satisfies \(|\eta_\varepsilon'|\le4/\varepsilon\).  Multiplying
also by a fixed outer cut-off equal to \(1\) on a ball containing
\(E\cup B\), set
\[
   X_\varepsilon(x):=\eta_\varepsilon(|x|)X(x).
\]
Then \(X_\varepsilon\) is Lipschitz and compactly supported, so the
Gauss--Green formula for finite-perimeter sets gives
\[
   \int_{\partial^*E}X_\varepsilon\cdot\nu_E\,d\Hn
      =\int_E {\rm div}\,X_\varepsilon\,dx,
   \qquad
   \int_{\partial B}X_\varepsilon\cdot\nu_B\,d\Hn
      =\int_B {\rm div}\,X_\varepsilon\,dx.
\]
The boundary integral over \(\partial B\) is zero because \(w(\rho)=0\).
Letting \(\varepsilon\downarrow0\) is legitimate: \(X_\varepsilon\to X\)
boundedly on \(\partial^*E\), \({\rm div}\,X\in L^1_{\rm loc}\) for
\(n\ge2\), and the cut-off derivative contributes at most
\[
   \frac{C}{\varepsilon}
   |(E\cup B)\cap\{\varepsilon/2<|x|<\varepsilon\}|
   \le C\varepsilon^{n-1}\to0.
\]
Hence
\[
   \int_{\partial^*E}X\cdot\nu_E\,d\Hn
   =
   \int_{E\setminus B}{\rm div}\,X\,dx
   -\int_{B\setminus E}{\rm div}\,X\,dx.
\]
For \(r=|x|>\rho\),
\[
   {\rm div}\,X=\frac{n}{\rho}-\frac{n-1}{r}
      \le \frac{n}{\rho}\le\frac{n}{\rstar}.
\]
For \(0<r<\rho\),
\[
   {\rm div}\,X=\frac{n-1}{r}-\frac{n}{\rho},
   \qquad
   (-{\rm div}\,X)_+\le\frac{n}{\rho}\le\frac{n}{\rstar}.
\]
Therefore
\[
   \int_{\partial^*E}X\cdot\nu_E\,d\Hn
   \le \frac{n}{\rstar}\bigl(|E\setminus B|+|B\setminus E|\bigr)
   =\frac{n}{\rstar}|E\Delta B|.
\]
Combining the two estimates proves the theorem.
\end{proof}

\begin{corollary}[Radial and homothetic traces on good levels]
\label{cor:fjbridge-homothetic-trace}
With the measurable same-centre field from
Proposition~\ref{prop:fjbridge-measurable-selection}, for a.e.
\(\rho\in G_I\),
\[
   T_\rho:=
   \int_{\partial^*E_\rho}
      \left|\frac{|x-z_\rho|}{\rho}-1\right|\,d\Hn(x)
   \le C(n,R,\rstar)\mathcal I(\rho)^{1/2},
\]
and, if
\[
   H_{z_\rho,\rho}(x):=
      \frac{(x-z_\rho)\cdot\nu_\rho(x)}{\rho},
   \qquad
   \bar V_\rho:=\frac{P(B_\rho)}{P(E_\rho)},
\]
then
\[
   \left(
      \int_{\partial^*E_\rho}
         |H_{z_\rho,\rho}-\bar V_\rho|\,d\Hn
   \right)^2
   \le C(n,R,\rstar)\mathcal I(\rho).
\]
\end{corollary}

\begin{proof}
Theorem~\ref{thm:fjbridge-radial-trace} and
Proposition~\ref{prop:fjbridge-good-level} give
\[
   T_\rho\le
   C\left(\mathcal I(\rho)^{1/2}+\mathcal I(\rho)\right).
\]
Since \(\mathcal I(\rho)\le\eta_I\le1\) on \(G_I\), this gives the
first estimate.

Writing \(e_\rho=(x-z_\rho)/|x-z_\rho|\), one has
\[
   |H_{z_\rho,\rho}-1|
   \le
   \left|\frac{|x-z_\rho|}{\rho}-1\right|
   +|e_\rho\cdot\nu_\rho-1|
   =
   \left|\frac{|x-z_\rho|}{\rho}-1\right|
   +\frac12|e_\rho-\nu_\rho|^2.
\]
After integration this is bounded by \(C\mathcal I(\rho)^{1/2}\).  Also
\[
   P(E_\rho)|1-\bar V_\rho|
   =P(E_\rho)-P(B_\rho)
   \le C(n,\rstar)\mathcal I(\rho)
   \le C(n,\rstar)\mathcal I(\rho)^{1/2}.
\]
The triangle inequality and squaring yield the homothetic trace bound.
\end{proof}

\begin{remark}[Import point for \textsc{WeightedMetricTrace}]
\label{rem:fjbridge-import-point}
The intended replacement for the current centre package in
\texttt{WIP\_WeightedMetricTrace.tex} is:
Proposition~\ref{prop:fjbridge-measurable-selection} for the Borel
same-centre field, Theorem~\ref{thm:fjbridge-radial-trace} for the
finite-perimeter oriented divergence trace, and
Corollary~\ref{cor:fjbridge-homothetic-trace} for the good-level
homothetic trace estimate used in the metric first-variation step.
\end{remark}
